%! AP Statistics — Unit 1, Lesson 1.5 — Group Activity ANSWER KEY
\documentclass[10pt]{article}
\usepackage{apstats-article}
\usepackage{apstats-key}

%=============================================================================
\begin{document}

\pageheader{Unit 1, Lesson 1.5}{Group Activity --- Answer Key}
\begin{tabularx}{\linewidth}{X X X}
Name:\;\ans{Key} & Partner:\; & Period:\;
\end{tabularx}

\vspace{0.1in}

% ── PART 1 KEY ───────────────────────────────────────────────────────────────
{\large\bfseries\color{navy} Part 1 \quad Read a Dotplot}

\vspace{0.1in}

\begin{center}
\begin{tikzpicture}[scale=1.05]
  \draw[thick] (3.2,0) -- (9.8,0);
  \foreach \x/\lbl in {4/4, 5/5, 6/6, 7/7, 8/8, 9/9} {
    \draw (\x,-0.12) -- (\x,0.12);
    \node[below, font=\small] at (\x,-0.12) {\lbl};
  }
  \node[below, font=\small\itshape] at (6.5,-0.52) {Height (cm)};
  \filldraw[navy] (4, 0.38) circle (0.14);
  \foreach \k in {1,...,3} \filldraw[navy] (5, \k*0.38) circle (0.14);
  \foreach \k in {1,...,4} \filldraw[navy] (6, \k*0.38) circle (0.14);
  \foreach \k in {1,...,4} \filldraw[navy] (7, \k*0.38) circle (0.14);
  \foreach \k in {1,...,2} \filldraw[navy] (8, \k*0.38) circle (0.14);
  \filldraw[navy] (9, 0.38) circle (0.14);
\end{tikzpicture}
\end{center}

\vspace{0.1in}

\begin{enumerate}[leftmargin=1.5em, itemsep=6pt]

  \item \ans{$4 + 4 = 8$ seedlings.}
        \ans{$8/15 \approx 0.533 = 53.3\%$.}

  \item \ans{Range $= 9 - 4 = 5$ cm.}
        \ansline{There are no gaps. There is a cluster in the 6--7 cm range
        (8 of the 15 seedlings), with fewer observations toward the extremes.}

  \item \ansline{Yes, the distribution is roughly symmetric. The dotplot has a
        peak at 6 and 7 (4 dots each), with 3 dots at 5 and 2 at 8 on either
        side, and 1 dot at each extreme (4 and 9).}

  \item \ansline{Yes, a dotplot is reasonable for $n = 15$ because it preserves
        every individual height. With 200 seedlings, dots would stack so high
        that the graph would be unwieldy and hard to read --- a histogram would
        be a better choice.}

\end{enumerate}

\vspace{0.1in}

% ── PART 2 KEY ───────────────────────────────────────────────────────────────
{\large\bfseries\color{navy} Part 2 \quad Back-to-Back Stemplot --- Key}

\vspace{0.08in}

\begin{center}
\renewcommand{\arraystretch}{1.8}
\begin{tabular}{r c l}
\textbf{Period 1} & \textbf{Stem} & \textbf{Period 2} \\
\hline
\ans{8} & 6 & \ans{5\; 9} \\
\ans{9\; 6\; 2\; 1} & 7 & \ans{2\; 5\; 7\; 8} \\
\ans{8\; 5\; 3} & 8 & \ans{2\; 4\; 8} \\
\ans{4\; 1} & 9 & \ans{1} \\
\end{tabular}
\end{center}

\begin{teachernote}
Period 1 leaves are written right-to-left (closest to stem first).
Period 2 leaves are written left-to-right. Both sets must be in ascending order
(reading outward from the stem).
\end{teachernote}

\begin{enumerate}[leftmargin=1.5em, itemsep=6pt]

  \item \ans{Period 1 ($n = 10$): median $=$ average of 5th and 6th values
        $= (79 + 83)/2 = 81$.}\\
        \ans{Period 2 ($n = 10$): median $=$ average of 5th and 6th values
        $= (77 + 78)/2 = 77.5$.}

  \item \ansline{Period 1 has a higher median (81 vs.\ 77.5). Period 2 has more
        spread: its range is $91 - 65 = 26$, compared to $94 - 68 = 26$ for Period 1
        --- the ranges are equal, but Period 2 has more scores in the 60s, pulling
        it lower.}

  \begin{teachernote}
  Range is 26 for both classes; accept this observation. Students might also
  note that Period 2 has 2 values in the 60s while Period 1 has only 1, giving
  more lower-tail spread. Require specific numerical evidence.
  \end{teachernote}

  \item \ansline{Period 1 has 5 scores at or above 80 (83, 85, 88, 91, 94),
        which is 50\% of its class. Period 2 also has 4 scores at or above 80
        (82, 84, 88, 91), which is 40\%. So Period 1 does have a slight edge in
        high scores, making the claim partially correct --- but Period 1's median
        (81) is only modestly higher than Period 2's (77.5), so the difference is
        not dramatic.}

\end{enumerate}

\vspace{0.1in}

% ── PART 3 KEY ───────────────────────────────────────────────────────────────
{\large\bfseries\color{navy} Part 3 \quad Histogram Interpretation --- Key}

\vspace{0.08in}

\renewcommand{\arraystretch}{1.6}
\begin{tabularx}{\linewidth}{@{} C{0.22\linewidth} C{0.18\linewidth} C{0.22\linewidth} C{0.26\linewidth} @{}}
\rowcolor{sky}
\textbf{Bin} & \textbf{Freq.} & \textbf{Rel.\ Freq.} & \textbf{Cum.\ Freq.} \\
\midrule
$[150, 160)$ & 3  & \ans{0.06} & \ans{3}  \\
$[160, 165)$ & 8  & \ans{0.16} & \ans{11} \\
$[165, 170)$ & 15 & \ans{0.30} & \ans{26} \\
$[170, 175)$ & 14 & \ans{0.28} & \ans{40} \\
$[175, 180)$ & 7  & \ans{0.14} & \ans{47} \\
$[180, 190)$ & 3  & \ans{0.06} & \ans{50} \\
\midrule
\rowcolor{sky}\textbf{Total} & 50 & \ans{1.00} & \ans{50} \\
\end{tabularx}

\vspace{0.1in}

\begin{enumerate}[leftmargin=1.5em, itemsep=6pt]

  \item (Answers above.)

  \item \ans{The $[165, 170)$ bin has the greatest frequency (15 adults, or 30\%).}

  \item \ansline{Because the bins are different widths, a taller bar does not
        necessarily mean a higher density of observations per unit of measurement.
        The $[150,160)$ bar spans 10 cm but the $[160,165)$ bar spans only 5 cm,
        so comparing their heights directly is misleading --- the wider bar appears
        taller even if fewer observations fall per centimeter.}

  \begin{teachernote}
  This is a preview of density histograms (optional). The key point is that equal
  bin widths are required for the bar height to represent frequency fairly.
  \end{teachernote}

  \item \ans{$(14 + 7 + 3)/50 = 24/50 = 0.48 = 48\%$ of adults are at least 170 cm tall.}

\end{enumerate}

\vspace{0.1in}

% ── PART 4 KEY ───────────────────────────────────────────────────────────────
{\large\bfseries\color{navy} Part 4 \quad Reflection --- Key}

\vspace{0.08in}

\begin{tcolorbox}[colback=goldbg, colframe=goldacc, boxrule=0.8pt, arc=2mm,
    left=3mm, right=3mm, top=2mm, bottom=2mm]
\small\textbf{Sample response:}\\[6pt]
\ans{The most important trade-off is dataset size vs.\ preserving individual values.
For the Part 1 seedling data ($n=15$), a dotplot was ideal because we could see
each height. For the Part 3 adult heights ($n=50$), a histogram grouped data into
bins and sacrificed individual values --- but it made the overall shape clearer and
would scale to 500 observations without becoming unreadable.}
\end{tcolorbox}

\end{document}
